% Generated draft: requires mathematical and source review.
\documentclass[11pt]{article}
\usepackage[margin=25mm]{geometry}
\usepackage{fontspec}
\setmainfont{Latin Modern Roman}
\usepackage{amsmath,amssymb,mathtools,unicode-math}
\setmathfont{Latin Modern Math}
\usepackage{enumitem,needspace,xurl,hyperref}
\setlength{\parindent}{0pt}
\setlength{\parskip}{5pt}
\setlength{\emergencystretch}{3em}
\newcommand{\N}{\mathbb{N}}
\newcommand{\Z}{\mathbb{Z}}
\newcommand{\Q}{\mathbb{Q}}
\newcommand{\R}{\mathbb{R}}
\newcommand{\C}{\mathbb{C}}
\newcommand{\F}{\mathbb{F}}
\newcommand{\A}{\mathbb{A}}
\newcommand{\PP}{\mathbb{P}}
\newcommand{\E}{\mathbb{E}}
\newcommand{\eps}{\varepsilon}
\newcommand{\dd}{\,\mathrm{d}}
\begin{document}
\section*{UnsolvedMath ID: 9700037}
\textbf{Source problem code:} catalogue entry
\subsection{Definitions and mathematical statement}
Let \(\mathcal P=\{C_i\}\) denote Aldous's source-specific self-similar random partition of \(\R^2\), with designated capital points \(c_i\). Its construction and probability law are not supplied. The task has several parts:
\begin{enumerate}
\item Determine topological properties of the regions and their boundaries, in particular whether boundary fractal dimension is greater than one. The supplied summary uses Hausdorff dimension, suggesting \(\dim_{\mathrm H}\partial C_i>1\), but the excerpt does not state a typical-region versus every-region quantifier or a dimension convention.
\item Determine the region area distribution. A region chosen by a typical capital and the region containing a typical location in the plane need not have the same law; sampling and scale normalization must be specified.
\item Determine the structural law of the capital-city adjacency network: vertices are the capitals and edges correspond to the source's region-adjacency relation. Is the abstract graph planar? Does its intended geometric drawing have crossings? Abstract planarity and absence of crossings in the straight-segment drawing between capitals are distinct questions.
\item Determine whether analogues of the Lewis and Aboav--Weaire laws hold. Ordinarily these concern area versus number of neighbors and neighbor-degree correlations, respectively. If \(A\) is a sampled region's area, \(D\) its degree, and \(M\) the average degree of its neighbors, familiar candidate relationships concern \(\E[A\mid D=d]\) as a linear function of \(d\), and \(\E[M\mid D=d]\) as a constant plus a reciprocal function of \(d\). The source gives no exact formulas, coefficients, error criterion, or sampling convention for its requested analogues.
\end{enumerate}
The adjacency rule, existence and finiteness of areas/degrees, dimension quantifiers, self-similar scale and law definitions require the omitted model. None is supplied by replacing the partition with a standard tessellation.

The registry reports that boundary Lebesgue measure zero and a strict Hausdorff dimension bound between one and two have been proved in related cited work, while the other laws remain open. These are supplied progress claims whose exact model and quantifiers must be checked, not new results asserted here.
\subsection{Short English statement}
Analyze the boundaries, region-area distribution, and capital adjacency graph of Aldous's particular random planar partition, including crossings and analogues of the Lewis and Aboav--Weaire correlations. The model and sampling conventions are missing; reported boundary results remain to be checked.
\subsection{Sources}
\begin{itemize}
\item Local UnsolvedMath registry, permanent ID 9700037, source code catalogue entry; David Aldous; no proposal year supplied; update 2026-08-06.
\item Aldous, \emph{Open problems}, partition.html: \url{https://www.stat.berkeley.edu/~aldous/Research/OP/index.html}. Registry access date 2026-07-29; HTML text with linked primary-PDF text. Provenance reports 23 active pages, with delegated explicit subproblems separated and compound bullets retained. Present status and redistribution terms require review.
\item The supplied summary cites Preater (2009) for zero boundary Lebesgue measure via a voter-model immigration analogue, and Basdevant--Blanc--Curien--Singh, \emph{ALEA} (2024), for strict Hausdorff dimension bounds between one and two. Full bibliographic entries, model identifications and exact quantifiers are not supplied.
\end{itemize}
\end{document}
